\documentclass[output=paper,colorlinks,citecolor=brown]{langscibook}
\ChapterDOI{10.5281/zenodo.20541866}


\author{Kristian Roncero\orcid{}\affiliation{Ludwig-Maximilians-Universität München}}

\title{West Polesian future tense constructions switch roles}

\abstract{In this chapter I assess West Polesian (East Slavic) future tense constructions from historical and typological perspectives, focusing on the de-obligative (xoʧu ɾobɪtɪ ‘I will (need to) do’) and de-volitive futures (majusʲ(a) ɾobɪtɪ ‘I am planning to do’), which have had the most unusual evolution of all. Against the typological expectations, the de-volitive has synchronically an obligational nuance, whereas the de-obligative expresses an intention. 

First, I describe their morphological structure and how speakers use them. Then I show their cross-linguistic recurrence and  presence in Slavic and other European languages. I asses the differences and similarities between the West Polesian constructions and their analogue forms or cognates and I map their distribution in Europe. In the light of this data, I try to address two closely related questions: (i) How and at which point may have this semantic change happened? (ii) Is the semantic mutation of these a side-effect of having a particularly large inventory of future tense constructions (which could have pushed some to specialize) in line with the Theory of Niches?
}

%move the following commands to the "local..." files of the master project when integrating this chapter
% \usepackage{tabularx}
% \usepackage{langsci-optional}
% \usepackage{langsci-gb4e}
% \bibliography{localbibliography}
% \usepackage{multirow}
% \newcommand{\orcid}[1]{}
% \let\eachwordone=\itshape




\IfFileExists{../localcommands.tex}{
   \addbibresource{../localbibliography.bib}
   \usepackage{tabularx,multicol}
\usepackage{url}
\urlstyle{same}

\usepackage[english]{babel}
 
\usepackage{langsci-optional}
\usepackage{langsci-lgr}
\usepackage{langsci-gb4e}

\usepackage{langsci-branding}

%\usepackage{tabularx}

% elena
\usepackage{graphicx} %package to manage images
\usepackage{paralist}

% sarah & lena
\usepackage[justification=centering]{caption}
\usepackage{rotating}

% robin
\usepackage{lscape}
\usepackage{tikz}
\usetikzlibrary{tikzmark}

% dominika
\usepackage{multirow}
\usepackage{hanging}

% olaf & dylan
\usepackage{subcaption}

% svetlana
\usepackage{verbatim}
\usepackage{enumitem}

% wiemer
%\usepackage[markup=underlined]{changes}
%\usepackage{hyphenat}

% wir
\usepackage{longtable}

% for crossref in the introduction jiabin
\usepackage{xr}
   % \newcommand*{\orcid}{}
%

\makeatletter
\let\thetitle\@title
\let\theauthor\@author
\makeatother

\newcommand{\togglepaper}[1][0]{
   \bibliography{../localbibliography}
   \papernote{\scriptsize\normalfont
     \theauthor.
     \titleTemp.
     To appear in:
     E. Di Tor \& Herr Rausgeberin (ed.).
     Booktitle in localcommands.tex.
     Berlin: Language Science Press. [preliminary page numbering]
   }
   \pagenumbering{roman}
   \setcounter{chapter}{#1}
   \addtocounter{chapter}{-1}
}
%
\newbool{bookcompile}
\booltrue{bookcompile}
\newcommand{\bookorchapter}[2]{\ifbool{bookcompile}{#1}{#2}}
%
% %\newcommand{\orcid}[1]{}
\graphicspath{ {./figures/} }
%
\let\eachwordone=\itshape


% Seongha
% \newcommand{\koreex}[1]{{\fontspec{Baekmuk Batang}{#1}}}

% Kutida
\newfontfamily\thaifont{Noto Sans Thai} 
\newcommand{\textthai}[1]{{\thaifont #1}} % für Thai


\newfontfamily\chinesefont{Noto Sans CJK SC}
\newcommand{\textchinese}[1]{{\chinesefont #1}} % Für Chinesisch

% wiemer
\def\hyph{-\penalty0\hskip0pt\relax}
% \definechangesauthor[color=NavyBlue]{SH}



\newcolumntype{K}{>{\raggedright\arraybackslash}p{0.6\textwidth}}  % Right column with 0.7 width


% \newfontfamily\krn[Scale=MatchLowercase,BoldFont=SourceHanSerif-Bold.otf]{SourceHanSerif-Regular.otf}
% \AdditionalFontImprint{Source Han Serif KO}
% \XeTeXlinebreaklocale 'ko'
% \renewcommand{\krn}{}



% Rhee
\newfontfamily\koreanfont{Noto Sans CJK KR} 
\newcommand{\textkorean}[1]{{\koreanfont #1}} % für Koreanisch

\renewcommand{\REF}[2][]{(\ref{#2}#1)}
   %% hyphenation points for line breaks
%% Normally, automatic hyphenation in LaTeX is very good
%% If a word is mis-hyphenated, add it to this file
%%
%% add information to TeX file before \begin{document} with:
%% %% hyphenation points for line breaks
%% Normally, automatic hyphenation in LaTeX is very good
%% If a word is mis-hyphenated, add it to this file
%%
%% add information to TeX file before \begin{document} with:
%% %% hyphenation points for line breaks
%% Normally, automatic hyphenation in LaTeX is very good
%% If a word is mis-hyphenated, add it to this file
%%
%% add information to TeX file before \begin{document} with:
%% \include{localhyphenation}

\hyphenation{
    con-tri-bute
    par-a-digm
    pe-ri-phra-ses
    achieve-ment
    accomp-lish-ment
    We-li-nchen-kang-ci-kok
    cross-ling-u-is-tic
    rec-ord-s
    trun-ca-ted
    de-ri-va-tion-al
    stem-de-ri-va-tion-al
    pre-sent
    in-flu-en-ced
    sys-tem
    ques-tions
    Bor-kovs-kij
    Da-ny-len-ko
    e-qui-va-lents
    re-cog-nized
    In-ter-fe-renz-schicht
    Aus-ei-nan-der-se-tz-ung
    Früh-neu-hoch-deutsch-kor-pus
    Ak-tions-art
    Smir-no-va
    Whi-le
    Ar-chive
    bi-ginnan
    Rei-chen-bach
    ger-ma-nis-ti-sche
    Re-fe-renz-kor-pus
    le-xi-ka-lisch-en
    Alt-hoch-deutsche
    Li-te-ra-tur
    dia-chron
    Dia-lek-to-lo-gie
}

\hyphenation{
    con-tri-bute
    par-a-digm
    pe-ri-phra-ses
    achieve-ment
    accomp-lish-ment
    We-li-nchen-kang-ci-kok
    cross-ling-u-is-tic
    rec-ord-s
    trun-ca-ted
    de-ri-va-tion-al
    stem-de-ri-va-tion-al
    pre-sent
    in-flu-en-ced
    sys-tem
    ques-tions
    Bor-kovs-kij
    Da-ny-len-ko
    e-qui-va-lents
    re-cog-nized
    In-ter-fe-renz-schicht
    Aus-ei-nan-der-se-tz-ung
    Früh-neu-hoch-deutsch-kor-pus
    Ak-tions-art
    Smir-no-va
    Whi-le
    Ar-chive
    bi-ginnan
    Rei-chen-bach
    ger-ma-nis-ti-sche
    Re-fe-renz-kor-pus
    le-xi-ka-lisch-en
    Alt-hoch-deutsche
    Li-te-ra-tur
    dia-chron
    Dia-lek-to-lo-gie
}

\hyphenation{
    con-tri-bute
    par-a-digm
    pe-ri-phra-ses
    achieve-ment
    accomp-lish-ment
    We-li-nchen-kang-ci-kok
    cross-ling-u-is-tic
    rec-ord-s
    trun-ca-ted
    de-ri-va-tion-al
    stem-de-ri-va-tion-al
    pre-sent
    in-flu-en-ced
    sys-tem
    ques-tions
    Bor-kovs-kij
    Da-ny-len-ko
    e-qui-va-lents
    re-cog-nized
    In-ter-fe-renz-schicht
    Aus-ei-nan-der-se-tz-ung
    Früh-neu-hoch-deutsch-kor-pus
    Ak-tions-art
    Smir-no-va
    Whi-le
    Ar-chive
    bi-ginnan
    Rei-chen-bach
    ger-ma-nis-ti-sche
    Re-fe-renz-kor-pus
    le-xi-ka-lisch-en
    Alt-hoch-deutsche
    Li-te-ra-tur
    dia-chron
    Dia-lek-to-lo-gie
}
   \boolfalse{bookcompile}
   \togglepaper[12]%%chapternumber
}{}

\begin{document}
\maketitle

\section{Introduction} \label{sec:roncero:1}
\subsection{A Quixotesque plot-twist} \label{sec:roncero:1.1}

By the end of Miguel de Cervantes’ “Don Quixote”, there is an unexpected plot-twist. Out of contact with each other, Don Quixote, who had been always the insane character that suffered from numerous hallucinations, dies sane; whilst his loyal squire Sancho, who used to be the voice of reason, turns insane. In this chapter, I describe an analogous situation concerning the course of two future tense constructions in West Polesian (East Slavic). Against the typological expectations, the de-obligative construction (i.e. historically based on a ‘have-to’ verb) has ended up acquiring a volitional/intentional nuance \textit{majusj(a) robɪtɪ} ‘I will/intend to do’; whilst the de-volitive construction (i.e. historically based on a ‘want/will’-type of verb) has turned obligative \textit{xoʧu robɪtɪ} ‘I will be in the need to do it’.

First, I start by presenting the main future tense constructions in West Polesian (\sectref{sec:roncero:1.2}), following the criteria set by \citet{Dahl2000b} to give an overview of the complexity of the set and its implications. I also discuss the sociolinguistic setting very briefly and provide some methodological remarks (\sectref{sec:roncero:1.3}). Second, I narrow the attention to the two future tense constructions under discussion, their morphological structure, and their uses (\sectref{sec:roncero:2}). Third, I look at the constructions from both an areal/typological and a diachronic perspective to assess whether they are loans from neighbors and how did they get into the system (\sectref{sec:roncero:3}). Finally, I present a summary of the findings and extract some conclusions (\sectref{sec:roncero:4}).

\subsection{Overview of future tense constructions in West Polesian}
\label{sec:roncero:1.2}
West Polesian is an East Slavic variety or language spoken in southwestern Belarus (Brest), northwestern Ukraine (Volynja) and a small region on Poland’s easternmost border (Podlasie). West Polesian is primarily a language of oral communication within the community, and speakers refuse to use it outside of their trusted community members, because it is highly stigmatized. The speech community has remained fairly isolated from the outside for centuries, given the swampy topography, which made travel both inside and outside of the community very difficult. As a result, culturally and linguistically they have remained more conservative in some respects than most of their Slavic neighbors (e.g. the presence of vocatives), but this has also set the ground for innovations (e.g. pronominal numerals; postnominal and dislocated possessors), some of which I discuss in this chapter. Conversely, for the past decades, the community has been surrounded by four, genetically very closely, yet highly standardized Slavic varieties: Standard Belarusian, Polish, Russian and Standard Ukrainian.

Narrowing down the focus on the peculiarities of West Polesian grammar, future tense constructions are certainly worth studying. In the light of \citegen{Dahl2000b} survey of future tense constructions in European languages, data from fieldwork in West Polesian future tense constructions fall into at least six of his future tense construction-families: 

\begin{enumerate}
    \item The \MakeUppercase{synthetic (imperfective) future} (cf. \citealt[319]{Dahl2000b}). It consists of the infinitive form of the verb plus a cliticised form of the verb ‘to have’ with \textbf{person/number} marking.

 \ea \label{ex:roncero:1}
     (Z10.1 00:20, own fieldwork)\\
     \gll ja ʃːjo nɪ zna-tɪ-mu skaza-tɪ\\
     1\textsc{sg.nom} \textsc{q.acc.sg} \textsc{neg} know-\textsc{inf-have.1sg} say-\textsc{inf}\\
     \glt ‘I won’t be able to tell [it] (lit. I will not know to tell).’ 
 \z

 \item The \MakeUppercase{de-obligative future} (\citeyear[323]{Dahl2000b}).\footnote{\citet[318--319]{Dahl2000b} calls this the “Balkan ‘have’ future”, although this is not the only de-obligative future-construction family in Europe. He previously mentions the “Romance inflectional future”, which is analogous to this form. However, for the sake of consistency and to avoid confusion with the rest of the terms here (i.e. not to confuse it with the inflectional or synthetic future, as in \textit{robɪtɪmu} ‘I will do’), I use a more generic term ‘de-obligative’.} It consists of a verb originally meaning ‘have to’ (in this case in the \textbf{reflexive} form) + a verb in the infinitive.
 \ea \label{ex:roncero:2}
     (T12.8.el 04:32, own fieldwork)\\
     \gll \textit{vona} \textit{maj-ɪt-sa} \textit{rodɪ-tɪ} \textit{ʋ} \textit{ɪjulj-e}\\
     3\textsc{sg.f.nom} have-\textsc{3sg-ref} deliver-\textsc{inf} in July-\textsc{loc.sg}\\
     \glt ‘She is due to (deliver in) July.’  
 \z

 \item The \MakeUppercase{de-volitive future} (cf. \citealt[323]{Dahl2000b}).\footnote{\citet[322]{Dahl2000b} mentions the existence of another de-volitive construction-family in Northern Europe (Germanic). Yet, the West Polesian de-volitive is genetically and typologically closer to the Balkan (i.e. South Slavic) family than to the Northern European one.} It consists of a verb originally meaning ‘to want’ + a verb in the infinitive.
 \ea \label{ex:roncero:3}
     (Tor1.36.pr 01:51, own fieldwork)\\
     \gll \textit{vona} \textit{xoʧ-e} \textit{kot-a} \textit{pokormɪ-tɪ}\\
     3\textsc{sg.f.nom} want-\textsc{3sg} cat-\textsc{acc.sg} feed.\textsc{perf}-\textsc{inf}\\
     \glt ‘She is going to feed the cat (lit. She’ll need to feed the cat).’  
 \z

\largerpage[2]

 \item The \MakeUppercase{‘become’-type} or \MakeUppercase{de-venitive future} (cf. \citealt[323--324]{Dahl2000b}). It consists of an auxiliary verb meaning ‘to become’ + a verb in infinitive.
 \ea \label{ex:roncero:4}
     (Tor1.53.el 03:28, own fieldwork)\\
     \gll \textit{jek} \textit{vona} \textit{stan-e} \textit{[xoroʃo]} \textit{ɦuʧɪ-tɪ-sja,} \textit{to} \textit{jɪjɪ} \textit{da-dutj} \textit{premɪj-u}\\
     \textsc{comp} 3\textsc{sg.nom.f} become.\textsc{prf-3sg} well study-\textsc{inf-refl} then 3\textsc{sg.dat.f} give.\textsc{prf-3pl} prize-\textsc{acc.sg}\\
     \glt ‘If/when she will start to study hard, she will get a prize.’  
 \z

 \item The Slavic \MakeUppercase{perfective future} (or \MakeUppercase{present}) (\citeyear[323, 326]{Dahl2000b}). It consists of a perfective verbal base, with \textbf{person/number}inflection as in the present tense.
 \ea \label{ex:roncero:5}
     (Z2.6 03:07, own fieldwork)\\
     \gll \textit{batjuʃk-a,} \textit{jak} \textit{ʒe} \textit{ʒ} \textit{ja,} \textit{[ɦet-oɦo],} \textit{sober-u} \textit{ɦet-ɪje} \textit{pir-ɪje?}\\
     priest-\textsc{voc} how then \textsc{part} 1\textsc{sg.nom} this-\textsc{gen.sg} gather.\textsc{prf-1sg} this-\textsc{acc.pl} feather-\textsc{acc.pl}\\
     \glt ‘Father (Orthodox priest), how will I gather all those feathers?’  
 \z

 \item The Slavic \MakeUppercase{copular constructions} (\citeyear[324]{Dahl2000b}), which consist of the auxiliary ‘to be’ in the future tense + a verb (in this case, in the infinitive).
 \ea \label{ex:roncero:6}
     (P2.5.pr 00:04, own fieldwork)\\
     \gll \textit{a} \textit{ʃːjo} \textit{mɪ} \textit{bud-em} \textit{robɪ-tɪ} \textit{na} \textit{pasku?} \textit{bud-em} \textit{praznova-tɪ}\\
     so \textsc{Q-acc.sg} 1\textsc{pl.nom} \textsc{be.fut-1pl} do-\textsc{inf} in Easter-\textsc{loc.sg}? \textsc{be.fut-1pl} celebrate-\textsc{inf}\\
     \glt ‘So what are we going to do for Easter? We are going to celebrate it.’  
 \z
 \end{enumerate}

On the one hand, this should not be surprising, as, according to the cross-linguistic surveys in \citet{Bybee1985} and \citet{Ultan1978}, future tense constructions are evolving frequently within verbal systems. Furthermore, West Polesian is on a highly diglossic sociolinguistic situation, so it is tempting to propose that the plethora of forms is poorly described code-switching or the result of extended language contact, which I try to challenge in this chapter. On the other hand, no European language appears to have more than three grammaticalized constructions for the future tense in \citegen{Dahl2000b} survey, which would make West Polesian a ‘record holder’ within Slavic languages, and most likely also in European languages. This represents an interesting situation from the point of view of typology, as well as historical linguistics, as this finding can assist in tracing areal influences.

\subsection{Methodological remarks} \label{sec:roncero:1.3}

Data for this chapter are the result of my own fieldwork in Belarusian Western Polesie (Brest region), between 2016--2017 as part of a larger morphological typology survey. Unfortunately, the pandemic, the migration crises with the border and the full-scale invasion of Ukraine have made impossible to go back to collect more data. I interviewed up to 50 speakers, however, for this parameter, I focused on free-texts and elicitation sessions with 15 speakers, as the constructions displayed are not present in all the varieties surveyed. The future tense has been by far the hardest parameter to elicit. Due to its delicate semantic nature (given that it is easily replaced by other forms) and relatively low frequency in the corpus, many of the examples provided here have been obtained through direct elicitation (most often using multiple choice questionnaires and grammaticality judgments, and, less frequently, using translations or rephrasing exercises). This is a common problem: \citet[89]{Ultan1978} also pointed out the lower saliency of the future in his cross-linguistic survey of the future tense. I have tried to be transparent about the source of all my examples. When the utterance was extracted from a free text, I have added the time stamp. In other cases, I have indicated whether it was prompted (pr) (e.g. by asking about a situation where they would be very likely to use the future tense); directly elicited (el) (e.g. by a translation task or grammaticality judgment); or overheard (ov) (and thus, not recorded) but outside of the working sessions (i.e. through participatory observation; e.g. when farming with them). 

Since West Polesian is primarily an oral variety and that there is hardly any reference of a standard, I have transcribed all the examples following my ongoing analysis of West Polesian phonology using IPA. Yet, I capitalize proper nouns, which I admit is a license I have taken (despite IPA conventions). 

Following the ethical standards in linguistic fieldwork, I collected consent from all participants before making the recordings. Given data protection and privacy regulations I have given pseudonyms to all participants involved. The corpus contains many recordings with sensitive data; therefore, following the common ethical guidelines, I have decided not to make them public. However, for reproducibility purposes, I invite people to email me to ask for access to the corpus. The letters of the code indicate the village where the recording was made. I am providing a list of the participants (with the names they chose to be known as) for those who opted for having their identity disclosed and get recognition for their valuable contribution to this project in the \hyperref[sec:roncero:app1]{Appendix}.

\section{Morphology and use} \label{sec:roncero:2}
\subsection{The de-obligative future} \label{sec:roncero:2.1}
\subsubsection{Morphological structure} \label{sec:roncero:2.1.1}
The ‘have’-type or de-obligative future construction in West Polesian is composed by the auxiliary verb \textit{matɪ} ‘to have’ (synchronically) with the reflexive suffix and the infinitive. The auxiliary bears person/number and probably also tense information, whereas the infinitive encodes the lexical meaning and aspect.

\ea \label{ex:roncero:7}
     \gll \textit{ja} \textit{ma-ju-sj} \textit{(po)sadɪ-tɪ} \textit{kartopl-i}\\
     1\textsc{sg.nom} have-\textsc{1sg-refl} plant.\textsc{prf-inf} potato-\textsc{acc.pl}\\
     \glt ‘I am going to plant potatoes.’  
\z


In the villages where I have worked I have only documented the reflexive form of the auxiliary (\textit{matɪsj(a) +} inf) ‘to have’. Nevertheless, when I asked my language assistants, they all recognized the non-reflexive form (\textit{matɪ} + inf) as grammatical. They pointed out that it is used in other villages, yet nobody managed to specify a village/variety where they use it. In any case, even though I primarily deal with the reflexive form here, the semantic base of the auxiliary is the verb ‘to have’ and so I include it in the group of de-obligative future constructions.

\largerpage[1]

Concerning the morphological structure, it is remarkable that \textit{matɪ} inflects for past tense (whereas other auxiliaries/synthetic futures do not) and indicates prospectivity in the past; as in \REF{ex:roncero:8}. To give some context, the speaker was narrating how the Nazis had taken hostage her dad and other men in the village. The soldiers were about to shoot the hostages, but someone that knew her dad intervened, so suddenly someone appears with a message from a superior telling the soldiers to stop the execution.

\ea \label{ex:roncero:8}
     (Tor1.25 05:31)\\
     \gll \textit{vosjɪm} \textit{ʧoloˈvɪk} \textit{ma-l-ɪ-sja} \textit{rostrɪlje-tɪ}\\
     eight.\textsc{nom/acc} person.\textsc{gradnm} have-\textsc{pst-pl-refl} shoot.\textsc{prf-inf}\\
     \glt ‘They were about to shoot eight people.’  
 \z

The auxiliary \textit{majusj} can be used on its own as a free or unbound auxiliary, like the copular auxiliary \textit{budu}. There are few examples of this in the corpus, but one instance is particularly worthy of mention. TL1’s daughter was supposed to water the potatoes one day, but she did not. While she was explaining to her mum why she had not been able to water them, her mother asked “\textit{a maiʃsja?}”, meaning ‘but are you going to water the potatoes?’


\subsubsection{Uses and connotations} \label{sec:roncero:2.1.2}

The de-obligative future tense construction has a strong sense of intentionality (in some cases it even retains some obligational flavor). After asking some speakers about the meaning of \textit{majusj} \textit{robɪtɪ} ‘I will do’ (de-obligative), some (e.g. T2) said that it has the same meaning as \textit{budu robɪtɪ} (copular), and thus that there is no difference between the two. Yet after several weeks some speakers (T5, T10 and T12 i.a.) pointed out that the auxiliary \textit{majusj} of the de-obligative is the equivalent of the Russian \textit{sobirajus´} ‘I intend to’.

\largerpage[1]

In most instances the \MakeUppercase{de-obligative future} is inappropriate for general predictions. However, in situations such as the discussion of someone’s pregnancy, this form appears the preferred one \REF[, 10a]{ex:roncero:9a}, rather than the inflectional future \REF[, 10b]{ex:roncero:9b}. Therefore, there may be also some lexicalization behind the choice of one structure over another.\footnote{Concerning (\REF{ex:roncero:9b}) and (\REF{ex:roncero:10b}) the \MakeUppercase{synthetic} form is generally good for both general predictions and intentions. In the context of discussing pregnancy and birth \textit{rodɪtɪme} implies a greater certainty than the \MakeUppercase{de-obligative} form \textit{maetsa rodɪtɪ}, according to my consultants. Given that people can expect, but not know for sure when they are going to give birth, the synthetic form sounded inappropriate to all three of my consultants. However, when it comes to sex of the baby, T13 argued that there can be more certainty about it (thanks to technology), thus, it is okay to speculate (using the de-obligative) about the sex of the baby, but it is also acceptable to assert it quite firmly (if we know the results from the doctor).}


\ea (T9, T12.el) \label{ex:roncero:9}
    \ea \label{ex:roncero:9a}
    \gll \textit{vona} \textit{ma-et-sa} \textit{rodɪ-tɪ} \textit{ʋ} \textit{ɪjunje} \\
    3\textsc{sg.nom.f} have-\textsc{3sg-refl} deliver-\textsc{inf} in June-\textsc{loc.sg}\\

    \ex[*]{
    \gll \textit{vona} \textit{rodɪ-tɪ-me} \textit{ʋ} \textit{ɪjun-je}\\
    3\textsc{sg.nom.f} deliver-\textsc{inf-have.3sg} in June-\textsc{loc.sg}\\
    \glt `She is due (to deliver) in June.'
    } \label{ex:roncero:9b}
    \z
\z

\ea \label{ex:roncero:10}
    (T9, T12.el)\\
    \ea \label{ex:roncero:10a}
    \gll \textit{vona} \textit{ma-et-sa} \textit{rodɪ-tɪ} \textit{doʧk-u}\\
    3\textsc{sg.nom.f} have-\textsc{3sg-refl} deliver-\textsc{inf} daughter-\textsc{acc.sg}\\

    
    \ex \label{ex:roncero:10b}
    \gll \textit{vona} \textit{rodɪ-tɪ-me} \textit{doʧk-u}\\
    3\textsc{sg.nom.f} deliver-\textsc{inf-have.3sg} daughter-\textsc{acc.sg}\\
    \glt ‘She will give birth to a girl.’
\z
\z

According to T6, the use of this form implies a high level of intentionality, with a sense of commitment “\textit{namjereno, ʋʒe zadumano}” ‘purposefully/intended, already planned’,\footnote{(T6.3 01:57).} whereas T9 emphasized that the de-obligative expresses more uncertainty than the perfective:

\begin{quote}
    “So if it’s tomorrow, I will say “tomorrow \textit{pojɪdu} ‘I will [already] go [\MakeUppercase{perfective}]’ to Drahičyn, or Brest, or Pinsk… I will go tomorrow”. But if I say \textit{majusja jɪxatɪ} ‘I’m due to go [\MakeUppercase{de-obligative}]’ it is unknown when \textit{pojɪdu }‘I will go [\MakeUppercase{perfective}]” (T9.9.el 02:32) [My translation].
\end{quote}

According to HL1’s comments, there is a correlation between animacy and the morphological strategy used. I proposed \REF{ex:roncero:11a}, where the object is a human (children), and she objected to it saying that “you would say\textit{ maju}, if you had an item, but with children, you say \textit{treba}\sim\textit{tra}\textit{ʔa}”.

\ea \label{ex:roncero:11}
    \ea {(HL1.el)\\
    \gll *\textbf{\textit{ma-ju-(sja)}} \textit{dɪt-ej} \textbf{\textit{poloʒɪ-tɪ}} \textit{spa-tɪ}\\
    have-\textsc{1pl-refl} child-\textsc{gen.pl} lay.\textsc{prf-inf} sleep-\textsc{inf}\\
    } \label{ex:roncero:11a}

    \ex (HL1.vol) \textrangle{}\textrangle \label{ex:roncero:11b}\\
    \gll \textit{mɪnjɪ} \textbf{\textit{treba}} \textit{dɪt-ej} \textbf{\textit{poloʒɪ-tɪ}} \textit{spa-tɪ}\\
    1\textsc{sg.dat} necessary child-\textsc{gen.pl} lay.\textsc{prf-inf} sleep-\textsc{inf}\\
    \glt ‘I need to/will put the children to sleep.’
\z
\z

\ea \label{ex:roncero:12}
    \ea (HL1.el) \label{ex:roncero:12a}\\
    \gll \textit{praz} \textit{pjet} \textit{ɦod-oʋ} \textbf{\textit{maj-ut-sja}} \textbf{\textit{dostroɪ-tɪ}} \textit{doroɦ-u}\\
    in five year-\textsc{gen.pl} have-\textsc{3pl-refl} build.\textsc{prf-inf} road-\textsc{acc.sg}\\
    \ex (HL1.vol) \textrangle{}\textrangle \label{ex:roncero:12b}\\
    \gll \textbf{\textit{sobiraj-ut-sja}} \textit{dostroɪ-tɪ}; \textit{objazan-ɪ} \textit{dostroɪ-tɪ}\\
    plan-\textsc{3pl-refl} build.\textsc{prf-inf} obliged-\textsc{nom.pl} build.\textsc{prf-inf}\\
    \glt ‘In five years’ time they will have to have finished building the road.’
\z
\z

% \largerpage[1]

\subsection{The de-volitive future} \label{sec:roncero:2.2}
\subsubsection{Morphological structure} \label{sec:roncero:2.2.1}
Some of the varieties surveyed in the West (of Belarusian Western Polesie) have a de-volitive future. In these varieties the verb \textit{xotitɪ} ‘to want’ has a deontic meaning, and it seems that \textit{majusj} (de-obligative) is close in meaning to this, as it was the first response I obtained whenever I tried eliciting sentences containing \textit{maju~/~majusj(a)}.

The de-volitive future tense is composed by an inflected form of the verb \textit{xotitɪ} ‘want’+ infinitive. Most commonly the verb \textit{xotitɪ} appears inflected in the present, which means that the topic-time is in the future; but less frequently the verb can appear in the past tense if the topic-time is in the past (i.e. future-in-the-past).\footnote{The future-in-the-past is less clear for the de-volitive. I have not been able to prove whether there are any tense restrictions, but I have a strong suspicion that there may be (unfortunately my main consultant, tor1, passed away during my research so i have not been able to confirm this). I have certainly identified utterances with the verb \textit{xotitɪ} in the past tense + infinitive in the corpus, although in most of them there were signs that it only had a ‘plain’ volitive meaning (i.e. no temporal or prospective value).} As with other periphrastic forms, the auxiliary encodes person, number and probably also tense information; whereas the infinitive bears the biggest semantic load (of the verb). In principle, both imperfective and perfective infinitives can appear with \textit{xotitɪ}. Nevertheless, when I suggested periphrases with \textit{xotitɪ} (among other constructions) with an infinitive to my main informant, Tor1, she rarely picked the imperfective infinitives with\textit{ xotitɪ}.

\ea \label{ex:roncero:13}
    (Tor1.el)\\
    \gll \textit{zaʋtra} \textbf{\textit{xoʧ-u}} \textbf{\textit{jɪxa-tɪ}} \textit{ʋ} \textit{boljnɪʦ-u}, \textit{bo} \textit{treba} \textit{operaʦɪj-u} \textit{zrobɪ-tɪ}\\
    tomorrow want-\textsc{1sg} go-\textsc{inf} to hospital-\textsc{acc.sg} as need operation-\textsc{acc.sg} do-\textsc{inf}\\
    \glt ‘Tomorrow I will have to go to the hospital to have an operation.’
\z

\ea \label{ex:roncero:14}
    (Tor1.el)\\
    \gll \textit{ja} \textit{xoʧ-u} \textit{zaruba-tɪ} \textit{pɪʋnj-a}, \textit{bo} \textit{vɪn} \textit{v} \textit{mena} \textit{klu-je} \textit{i} \textit{moʒe} \textit{i} \textit{popad-e} \textit{do} \textit{dɪt-aj}\\
    1\textsc{sg.nom} want-\textsc{1sg} kill.\textsc{prf-inf} rooster-\textsc{acc.sg} as 3\textsc{sg.nom.m} in 1\textsc{sg.acc} bite-\textsc{3sg} and may and attack.\textsc{prf-2sg} to child-\textsc{gen.pl}\\
    \glt ‘I will need to kill the rooster as he keeps biting, and he may attack children.’
\z

% \largerpage[1]

As elsewhere in East Slavic, there is also a verb \textit{xotitɪ\textsuperscript{2}} which means ‘to want’. It is hard to determine whether it is a completely different entry at this stage, but it is very clear that the de-volitive future (synchronically with a deontic value) originated from it.

\ea \label{ex:roncero:15}
    (Z10.12 00:30)\\
    \gll \textit{a} \textit{tɪpera} \textit{nɪkto} \textit{nɪ} \textit{xoʧ-e} \textit{za} \textit{ɦroʃ-ɪ} \textit{do} \textit{tibe} \textit{i-tɪ}\\
    and now nobody.\textsc{nom.sg} \textsc{neg} want-\textsc{3sg} for money-\textsc{acc.pl} to 2\textsc{sg.gen} go-\textsc{inf}\\
    \glt ‘So nowadays nobody wants to come [to help] you [even] for money.’
\z

\subsubsection{Uses and connotations} \label{sec:roncero:2.2.2}
The de-volitive future has very restricted uses as a future tense periphrasis. The verb \textit{xotitɪ} denotes a need, a type of obligation that is not (at least directly) imposed by external forces, such as the weather, or an institution like a bank.\footnote{Otherwise, other forms are used; e.g. \textit{ja kraj maju} (T6), \textit{ja dolʒɪn, mɪnjɪ musovo} or \textit{mɪnjɪ treba}. Tor1 was also aware of this restriction when I tried eliciting different forms (Tor1.45) “\textit{ja xoʧu, ɦato, ʃ:jo to takoje, znaʧit ʧastnɪ, ʃ:jo ja [no] xoʧu; ɦeto nado, treba platɪtɪ”} ‘ja xoʧu [I have to] is something private, but [with your bills] you must pay [them]’.} Note that in \REF{ex:roncero:13} the speaker insisted they did not want to go through an operating theatre, it was because they needed the procedure. If the verb \textit{xotitɪ} is used with an externally imposed obligation it loses its deontic meaning and only has (or preserves) a volitive meaning, which as a consequence sounds unacceptable to the speakers.

\ea \label{ex:roncero:16}
   (Tor1.el) [Context: A man has been outside the house for a while. He really has to get into the house quietly.]\\
    \gll vɪn xoʧ-e tɪxenko prɪ-tɪ bo djɪtɪ sp-jetj\\
    3\textsc{sg.nom} want-\textsc{3sg} silently enter-\textsc{inf} because child-\textsc{nom.pl} sleep-\textsc{3pl}\\
    \glt ‘He will need to get in silently because the children are sleeping.’
\z

\ea \label{ex:roncero:17}
\ea \label{ex:roncero:17a}
    (Tor1.el) [Context: There hasn't been any rain for a while.]\\
    \gll *ja xoʧ-u polɪ-tɪ oɦorod, bo vs-jo posoxn-e\\
    1\textsc{sg.nom} want-\textsc{1sg} water-\textsc{inf} garden-\textsc{acc.sg} because all-\textsc{nom.sg.n} dry.\textsc{prf-3sg}\\

\ex \label{ex:roncero:17b}
    (Tor1.vol) \textrangle{}\textrangle{} \\
    \gll treba (obɪzatilno) polɪ-tɪ oɦorod-a\\
    necessary obligatorily water-\textsc{inf} garden-\textsc{gen.sg}\\
    \glt ‘I (will) need to water the garden, or else everything will dry up.’
    \z
\z

Thus, the de-volitive form is closer to intentions than to predictions, but it still fails to fully adjust to any of these categories.

\section{Typological overview and diachrony} \label{sec:roncero:3}
\subsection{The de-obligative future} \label{sec:roncero:3.1}
\subsubsection{The de-obligative within East Salvic} \label{sec:roncero:3.1.1}
Before continuing further, \citet[323]{Dahl2000b} suggested a possible link between the (Ukrainian) synthetic future and the Balkan de-obligative form. Since West Polesian has both forms, I distinguish the de-obligative from the synthetic \textit{robɪtɪ-mu}, admitting that there may be some links between the two, but their grammaticalization paths have been very different.

Not all the contemporary Slavic languages have preserved a verb for ‘to have’ (particularly, in East Slavic), but there are obligative constructions in many, which use the verb ‘to have’ ((\textit{i)miec} or any of its variations); e.g. Polish \textbf{\textit{Mam}}\textit{ coś \textbf{robić} }‘I have to do something’. Neither have all the contemporary Slavic languages grammaticalised (or preserved) the verb ‘to have’ as a future tense auxiliary, but there are a few.

According to \citet[172]{Danylenko2011} and \citet[58]{Whaley2000} Rusyn or “Transcarpathian dialects” have a de-obligative future form, which is realised analytically, the auxiliary appearing before the infinitive. Nevertheless, \citet{Pugh2009} does not make any reference to this form in his grammar of Rusyn. \citet[61]{Whaley2000} briefly mentions that apparently there are de-obligative futures in North Russian dialects (more specifically she mentions the varieties spoken in the Velikij Ustjug, Grjazovec, Kadninkov and Čerepovec regions), although she does not provide any examples.

Belarusian (BEL) has the closest analogue to the West Polesian de-obligative. In both Standard Belarusian (particularly in Narkamaŭka) and West Polesian the verb (BEL) \textit{mec´/} (WP) \textit{matɪ} ‘to have’ is heavily eroded. \citet{Mazzitelli2011} ran a corpus survey on both \textit{mec´} and \textit{mecca} (the reflexive form) in Contemporary Standard Belarusian texts.\footnote{\citet[181]{Lomtev1956} says that, although very rarely, constructions with \textit{mec´} can be found in Contemporary Belarusian works and that they used to be the norm (and where widely used) in older stages of Belarusian. He claims that nowadays such constructions should be considered “dialectal”.}

Given that the verbs are rare \textit{per se}, not surprisingly, the results she obtained were that both forms are very infrequent, particularly the reflexive one: 0.05 per thousand words for \textit{mec´}; and 0.01 per thousand words for \textit{mecca} (\citeyear[180--181]{Mazzitelli2011}).\footnote{It must be also said that the reflexive form \textit{mecca} is has a very low frequency in both standards of Belarusian (\textit{Taraškevica }and\textit{Narkamaŭka), }whereas the contrast in the use of \textit{mec´} is more notorious between the standards (\cite[181]{Mazzitelli2011}).} According to \citet[200]{Mackevic1959} in Standard Belarusian, both \textit{mec´} and \textit{mecca }+ infinitive can be employed in the past and the present tenses.

Now, according to \citet[200]{Mackevic1959} and \citet[182]{Mazzitelli2011} in Belarusian, both \textit{mec´} and \textit{mecca} combine with perfective and imperfective infinitive forms and they do not impose semantic restrictions on the subject; i.e. they allow either animates or inanimates to be subjects of the verb. Thus, the West Polesian and Standard Belarusian de-obligative forms share many properties in common, and it is reasonable to infer a common origin.\footnote{I lack enough data to assert whether West Polesian imposes any animacy restriction on its subjects with de-obligatives. Nevertheless, I have observed \textbf{animacy} restrictions on the object of the de-obligatives (see \sectref{sec:roncero:2.1.2}, particularly \REF[--12]{ex:roncero:11}).}

Concerning the uses and meanings of both \textit{mec´} and \textit{mecca} +infinitive, \citet[200]{Mackevic1959} says that in Standard Belarusian they can have different modal meanings:

\begin{quote}
    “The modal meanings vary according to the form of the verb used: \textit{mec´} or the reflexive \textit{mecca}. The first one expresses a stronger meaning of obligation, need. […] Constructions with the reflexive form, \textit{mecca}, are characterized by a modal meaning of intention, disposition to undertake the action. […] The (modal) semantic nuances of these constructions are close to modal constructions of the type \textit{xacec´} [‘want’], \textit{žadac´} [‘wish’], \textit{dumac´} [‘think’], \textit{music´} [‘must’]” \citep[200]{Mackevic1959} [My translation].
\end{quote}

On the one hand, the Belarusian form \textit{mec´} lies between temporality and modality. Among the temporal notions it can express (mostly future), there are scheduled future (the most common), fatalistic future (especially when used in the past tense) and intentionality \citep[182, 190]{Mazzitelli2011}.

Very rarely, Belarusian \textit{mec´} can be used to express immediate future, but retaining a flavor of obligation, which for \citet[181]{Lomtev1956} and \citet[184]{Mazzitelli2011} disqualifies it as a fully grammaticalized (i.e. in their understanding, semantically neutral) future auxiliary. On the other hand, \citet[185]{Mazzitelli2011} says that the primary function of Contemporary Belarusian \textit{mecca} (i.e. the reflexive form) is “to express the notion of scheduled future in the past, usually with counterfactual interpretation […], and intention”.

\ea \label{ex:roncero:18}
    Standard Belarusian [See references in \citet[181]{Mazzitelli2011}]\\
    \gll Ja me-l-a-sja sta-c’ povar-am\\
    1\textsc{sg.nom} have-\textsc{pst-f.sg-refl} become-\textsc{refl} cook-\textsc{ins.sg}\\
    \glt ‘I should have become a cook.’
\z

Notwithstanding the semantic differences between the two forms, \citegen[180--181]{Mazzitelli2011} survey showed that \textit{mec´}is primarily used in the present tense, whereas \textit{mecca} (i.e. the reflexive form) almost exclusively in the past tense. That is to say, the main condition for one form or another seems to be \textbf{tense}. In contrast, it is true that West Polesian \textit{majusj(a)} + inf appears more often in the past tense (either as a counterfactual event or an intention). Yet it is still quite frequently employed in the present/future tense (i.e. as a proper future tense auxiliary), with an added meaning of intentionality. That is to say, the West Polesian reflexive auxiliary is not \textbf{tense}-bound, and displays more traits of a canonical or grammaticalized future tense construction.

 From a diachronic perspective, according to \citet[54]{Whaley2000} “[…] the fourteenth\hyp{}sixteenth centuries show widespread use of \textit{maju }plus the infinitive, a ‘have/take’-type future similar to that found in Old Church Slavonic as well as Rus´ianChurch Slavic”.\footnote{Conversely, the copular future was marginal at that time (\textit{ibid}).} By contrast, the reality I have found in \textit{The Belarusian Historical Dictionary} (henceforth, HSBM)\footnote{For its abbreviation in Belarusian (\textit{Histaryčny Sloŭnik Belaruskaj Movy}).} is quite different. I have searched for all the possible variants of the verb (\textit{mati; imati, meti}) with their respective reflexive counterparts in \textit{HSBM}. For the form \textit{mati} as ‘to own; to have’, there are only a few historical attestations of the verb and not much information is given. 

As for \textit{imati} and \textit{imatisja}, \textit{HSBM }very briefly says that both have been used as auxiliaries for the analytic future tense. The earliest instance of \textit{imati} that it documents is from 1438, where it says it has a modal/deontic nuance. And for \textit{imatisja}, the earliest token also comes from a fifteenth-century text (although no comments are made about any additional nuances).

\largerpage[-3]

Finally, the pair \textit{meti} and \textit{metisja} are the most prominent in the dictionary (HSBM vol. 18), and yet there is hardly any information available. Out of the many meanings for \textit{meti}, it says that it was possible to use it as an auxiliary for future tense denoting obligation or possibility. The earliest potential example of a de-obligative future that it provides is from a text dated from 1533. Yet, whilst the deontic reading is very clear, the temporal value is very questionable. The available material for \textit{metisja} is even more scarce. Only the 7\textsuperscript{th} entry describes very concisely that it has been used as a future tense auxiliary but denoting obligation. It only cites an example from 1579 \REF{ex:roncero:19}:

\ea \label{ex:roncero:19}
    Old Belarusian (AVK, 406 (XVI century) In: HSBM vol. 18 [my glossing and translation])\\
    \gll by-l-a u mene t-aja prijatelk-a … o t-om ne veda-l, esli by ee me-l-i-se svata-c’ ljud-i abo ne\\
    be-\textsc{pst-f.sg} in 1\textsc{sg.gen} that-\textsc{nom.sg.f} girlfriend-\textsc{nom.sg.f} {} about that-\textsc{loc.sg} no know-\textsc{pst.m} COMP PART 3\textsc{sg.acc.f} HAVE-\textsc{pst-pl-refl} propose-\textsc{inf} person-\textsc{nom.pl} or no\\
    \glt ‘I used to have a girlfriend… I didn’t know whether they were going to offer her in marriage or not.’
\z

\citet[186]{Mazzitelli2011} points out that the reflexive form was already present in Old Belarusian, and it already differed from its neighbours. According to \citegen{Mazzitelli2011} corpus study (\citeyear[192--193, 199]{Mazzitelli2011}), Old Belarusian \textit{mecca} was already grammaticalized as an auxiliary; and it was available for the future and the past tenses (as in West Polesian), the future being the least common. Yet, already by the nineteenth century it was almost restricted to the past tense (\textit{mec´} being used for the future, instead), expressing intention or mostly a counterfactual meaning \citep[192, 193, 199]{Mazzitelli2011}. The reflexive past tense form was preserved into the twentieth century in Belarusian. The use of \textit{mecca} did not decrease as significantly as the Belarusian \textit{mec´} in Soviet Belarus, but it remained very much marginal. 

\citet{Mazzitelli2011} presents various theories about the grammaticalization path of the reflexive de-obligative auxiliary. According to her, one possible explanation is that the auxiliary may have emerged from one of the meanings the reflexive \textit{mecca} had in Old Belarusian ‘strive for, going to’ (without an infinitive) \citep[198]{Mazzitelli2011}. However, the theory that \citet[199--200]{Mazzitelli2011} favors most is that Old Belarusian \textit{meti }and \textit{metisja }emerged almost in parallel. With time\textit{ meti}would have retained some modal value and added a value of ‘external obligation’, whereas \textit{metisja} would have developed into a more neutral auxiliary to replace the gaps of \textit{meti,} particularly in the past tense. In addition to this, Mazzitelli finds this last explanation particularly helpful, because it helps to explain why the reflexive forms are overwhelmingly more common than the non-reflexive ones in all periods of Belarusian when used with first and second persons.

In the light of this, we can say that West Polesian may have diverged from the rest of the East Slavic varieties already by the sixteenth century by preserving the reflexive form as a future tense auxiliary. So, in this respect, West Polesian is more conservative than the rest of the East Slavic clade.

\subsubsection{Balkan de-obligatives}
Bulgarian uses two main forms as a basis for the rest of the future forms: the form \textit{šte }(a de-volitive which I discuss in (\sectref{sec:roncero:3.2.1})\footnote{In Macedonian the cognate is \textit{ḱe} \citep[323]{Dahl2000b}.} for the positive forms; and for the negative forms the particles \textit{da }(“an invariable conjunctive particle”) and \textit{njama}, which actually is “the 3\textsuperscript{rd} person singular present form of the verb \textit{njamam} ‘have not’” \citep[209]{Kuteva1995}. Even though the Bulgarian/Macedonian verbal system is more complex than that of West Polesian and Belarusian, one of the closest similarities is the possibility of expressing a counterfactual event grammatically \citep[242]{SussexCubberley2006}. Nevertheless, besides the semantic basis and its possibility of expressing a counterfactual past event (for both imperfective and perfective verbs, see \citealt[209--212]{Kuteva1995}), the differences between the West Polesian and Bulgarian/Macedonian forms are greater than their similarities. First, Bulgarian polarity restrictions of do not apply to the West Polesian de-obligative. Second, Bulgarian has gone further along the grammaticalization path. As a result, Bulgarian \textit{njama} is a semantically neutral form that can express both intentions and predictions; whereas West Polesian \textit{majusj(a)} has a strong intentional connotation. And third, Bulgarian \textit{njama} does not currently inflect for \textbf{person} and does not have a reflexive form, whereas the West Polesian de-obligative is particularly salient (in the Slavic family, at least) for having a reflexive form.

Finally, \citet[323]{Dahl2000b} gives examples of de-obligatives in Gheg Albanian and Romanian. \citegen[164]{Danylenko2011} comments suggest that the closest form (to East Slavic de-obligatives) in the Balkans is the Greek \textit{μέλλω} which means ‘to intend’. However, \citet[164--167, 177]{Danylenko2011} also believes that the emergence of the Balkan de-obligative may have a Slavic origin, but it would not be related to the forms that emerged in Middle Belarusian and Ukrainian.

\subsubsection{Within the wider European context}

Outside of the Balkans, the two best known de-obligative ‘future construction families’ are English \textit{shall} (which earlier meant ‘to owe’, see \citealt[194]{Bybee1985}) and the Romance simple future; e.g. [French] \textit{aimer+ai }>\textit{aimerai} ‘I will love’. In comparison to West Polesian (and East Slavic), Romance de-obligatives are exclusively employed in the future tense. That is to say, they cannot be used as a future-in-the-past; because they tend to use the conditional for that, which is also related to auxiliary \textit{habere} (e.g. [French] \textit{aimerais}) or most often an imperfective de-andative periphrasis. I have summarised the areas in Europe where there is a de-obligative in (\figref{fig:roncero:1}). Totally unrelated to the Balkans or Slavic, Basque has also a de-obligative future, which is particularly used in familiar speech (as an alternative to the genitive-based future). The future reading is deduced from the context; cf. \REF{ex:roncero:20}--\REF{ex:roncero:21}:

\ea \label{ex:roncero:20}
    Basque (own knowledge)
    
    [Context: a friend asking another friend before ordering.]\\
    \gll Kafe-a eska-tu \textbf{behar} (al) du-zu?\\
    coffee-\textsc{det.abs.sg} ask-\textsc{inf} need \textsc{q.part} have.\textsc{3sg<2sg}\\
    \glt ‘Are you going to order coffee? (lit. Do you have to order coffee?)’
\z

\ea \label{ex:roncero:21}
    [Context: Doctor talking to a patient.]\\
    \gll Ur gehiago edan  \textbf{behar} du-zu.\\
    water.\textsc{abs} more drink.\textsc{inf} need have.\textsc{3sg<2sg}\\
    \glt ‘You need to drink more water.’
\z

\begin{figure}
    \centering
    \includegraphics[width=0.8\linewidth]{figures/roncero/figure1.pdf}
    \caption{Distribution of de-obligative future in Europe. Sources (\cite{BybeeDahl1989}, \cite{Dahl2000b}, \cite{Danylenko2011}, \cite{Mazzitelli2011}, \cite{Whaley2000}). (CC-BY-SA Kristian Roncero).}
    \label{fig:roncero:1}
\end{figure}

Outside Europe, \citet[194]{Bybee1985} also mentions Kru dialects as having a de-obligative future. Hence, there are de-obligative futures in different corners of Europe and the world, with little affiliation between them (if any). In fact,according to \citegen[194]{Bybee1985} survey, modals denoting ‘obligation’ “tend to develop into \textit{epistemic }markers” and “into \textit{futures}, as in the case of English \textit{shall}”. \citet[118]{Ultan1978} proposed that the link between modals and their common grammaticalization into future tense markers cross-linguistically may be due “to the relative uncertainty inherent in both future events and most of the categories subsumed under the general heading of modals”.

Thus, my hypothesis so far is that the East and South Slavic de-obligative future would have departed from a common material (the verb ‘to have’), that due to its semantics ‘naturally’ grammaticalized into a future tense construction. In this process, it is noteworthy that Belarusian and especially West Polesian have developed a de-obligative future tense auxiliary based on a reflexive, which is nowhere else found in Europe, and that it expresses intention rather than obligation.

\largerpage[2]

\subsection{The de-volitive future} \label{sec:roncero:3.2}
\subsubsection{Within the Slavic family} \label{sec:roncero:3.2.1}
There is evidence of the existence of de-volitive and de-obligative periphrases in older stages of East and Common Slavic (\cite{Danylenko2011}, \cite{Mackevic1959}, \cite{Whaley2000}). Thus, since this type of construction has a long tradition in Slavic, it is not surprising to find it in other Slavic varieties.

Starting from East Slavic, \citet{Mackevic1959} vaguely mentions the existence of periphrastic constructions with the verb \textit{xacec´} in Belarusian. However, he does not provide any examples, nor does he say in which dialects this form is used (certainly, it is not part of the current standard BEL, so he may be referring to West Polesian or Southwestern Belarusian, in general). According to \citet[154]{Danylenko2011}, Transcarpathian (i.e. Rusyn) and Eastern Ukrainian also have de-volitive periphrastic constructions based on the verb \textit{xotity}, although they have not fully grammaticalized the verb into a future tense auxiliary yet.\footnote{Unfortunately, \citet{Danylenko2011} does not provide more details or examples; nor does \citegen{Pugh2009} Rusyn grammar make any mention of it.} In this respect, West Polesian is probably not the only member of the East Slavic clade to have preserved a de-volitive form. Nonetheless, it has gone further in the path of grammaticalization than any other member and it has made the biggest switch in the semantic content of the base: from volition to obligation.\footnote{HSBM (vol.37) does not record any deontic meaning of the verb \textit{xotitɪ}. Levancèvič (p.c.) shared with me that other West Polesian sub-varieties (as in Malaryta and southern Byarozy) also use \textit{xotitɪ} with a deontic meaning. She considers it to be a specificity of West Polesian.}

Moving further south, Bosnian-Croatian-Montenegrin-Serbian (BCMS hereafter) forms its future tense using the inflected form of the verb \textit{hteti}\textit{ / htjeti} ‘want’ as an auxiliary, which is a cognate of West Polesian \textit{xotitɪ}. When it is in a positive sentence, the auxiliary can appear either as the whole form (\textit{hóču/hóčeš/}etc. + infinitive) or as a shortened form (\textit{ču/češ/}etc. + infinitive). However, when the verb is negated, only the neg \textit{ne} + the short form is possible (\textit{nè ču }or \textit{nèču /ne češ }or \textit{nečeš} \textit{/}etc. + infinitive; but *\textit{ ne hóču/}etc. + infinitive) (\cite[143--144]{Alexander2006}).

\citet[186, 189, 190]{Gvozdanovic1995} says that BCMS future tense (regardless of \textbf{aspect}) is used: “(a) for a future event, as in \REF{ex:roncero:22}; (b) for a necessary or desired event, as in \REF{ex:roncero:23}”. \citet[186]{Gvozdanovic1995} shows the following examples of BCMS futures:

\ea \label{ex:roncero:22}
    BCMS \citep[186]{Gvozdanovic1995}\\
    \gll Sutra ć-u ići u grad.\\
    tomorrow will.\textsc{prs-1sg} go.\textsc{inf} in/to city.\textsc{acc.sg}\\
    \glt ‘Tomorrow I shall go to the city.’
\z

\ea \label{ex:roncero:23}
    BCMS \citep[186]{Gvozdanovic1995}\\
    \gll Sada ć-eš ići po to.\\
    now will.\textsc{prs-2sg} go.\textsc{inf} for this.\textsc{acc.sg}\\
    \glt ‘Now you will go to fetch it.’
\z

As we can appreciate from \citegen[186]{Gvozdanovic1995} data and discussion, particularly \REF{ex:roncero:23}, this de-volitive construction shares with its West Polesian cognate part of the necessity meaning (which is a further innovation, given the etymology of the gram). Moreover, both constructions combine with perfective as well as imperfective infinitives, which is an uncommon feature with regards to other constructions in WP (and East Slavic).\footnote{According to \citet[186--187, 192]{Gvozdanovic1995} the BCMS and Slovene ‘future perfects’ can be used with either perfective or imperfective verbs (though they are based on the copular form).} However, the usage of this form is far more limited in West Polesian than in BCMS. BCMS allows the use of this form as a semantically neutral future form (i.e. for both intentions and predictions); whereas West Polesian de-volitive implies the existence of a need to be addressed. In exchange, West Polesian, having a morphologically richer repertoire of constructions, most frequently uses the copularform in such a context (i.e. a semantically neutral future or a prediction).

Bulgarian/Macedonian also has a future tense auxiliary which was etymologically a volitional verb, \textit{šta} ‘will/want’, (cognate and related of BCMS \textit{hteti/htjeti}).\footnote{According to Spencer (p.c.) in contemporary BCMS that verb has been replaced by \textit{iskam.}} Nonetheless, the auxiliary is almost exclusively used in positive sentences (with the vast majority of negatives it uses the de-obligative). The future (simple) is formed by the auxiliary \textit{šte}\textit{/ḱe}, etymologically the 3sg form of the verb \textit{šta} (although the auxiliary itself does not inflect for \textbf{person}), and the present tense form of the main verb \citep[209]{Kuteva1995}.


\ea \label{ex:roncero:24}
    Bulgarian \citep[209]{Kuteva1995}\\
    \gll Šte dadeš\\
    will/want(\textsc{invar}) give.\textsc{pfv.2sg.prs}\\
    \glt ‘You will give.’
\z

The Bulgarian de-volitive future form is versatile with regard to \textbf{tense }and is the base of the future simple; future perfect; future-in-the-past and future perfect-in-the-past forms \citep[209--212]{Kuteva1995}.\footnote{Note that in the future-in-the-past and the future-perfect-in-the-past the auxiliaries inflect for \textbf{person }as well as \textbf{tense}.} Even the future simple form (as in \REF{ex:roncero:24}), can express “an activity which will take place […] or will be taking place […] after the moment of speech” or “a repeated or habitual activity either in the present or in the past” \citep[209]{Kuteva1995}.\footnote{Spencer (p.c.) pointed out to me that English \textit{will} futures can also express generic habitual events or past events; e.g. “\textit{By 1811 Napoleon will become master of half of Europe; by 1815 he will be defeated and in exile}”. He also remarked that in Bulgarian renarrated or evidential mood you also encounter future tenses.} Conversely, the West Polesian de-volitive is very restricted in \textbf{tense} (in contrast with other West Polesian future-tense constructions also discussed here): it can only be used in the future tense (by inflecting the auxiliary in the present tense), as in any other form it functions as a plain volitive verb (see \sectref{sec:roncero:2.1}).

Elsewhere in the Balkans, \citet[323]{Dahl2000b} mentions further languages with de-volitives. On the one hand Modern Greek and Gheg Albanian which, in the same vein as Bulgarian/Macedonian, have an uninflectable form of the verb ‘want’ as the semantic base of their future tense auxiliary (\citealt[257]{Duchet1995}, \citealt[233]{Hedin1995}).\footnote{In Modern Greek (\textit{tha}) and Bulgarian/Macedonian (\textit{šte/ḱe}) the verb of the auxiliary has not survived as an independent volitional verb to our days.} And on the other hand Romanian, like BCMS, has an inflected auxiliary meaning ‘want’. Having said this, we should not be surprised to find more de-volitive future constructions in the world’s languages. Since volition and intentionality are semantically very closely related, volitive verbs often develop into future tense constructions (\cite{Bybee1985}, \cite{BybeeDahl1989}).

Finally, even though small, it is worth noting that Resian Slavic (RES) has also a copular, de-volitve and a de-obligative future constructions (\cite{Pila2022}), which is not only unique in South Slavic but it is also a rarity within Slavic. These constructions are not exactly interchangeable and all add some semantic nuances. In the case of the de-volitive \REF{ex:roncero:25}, it expresses more certainty than the copular form (be + participle) \REF{ex:roncero:26}, according to \citet[38]{Pila2022}. As for the de-obligative \REF[--28]{ex:roncero:27}, \citet[41]{Pila2022} states that it has more of a predictive function “the need of fulfilling an action, which depends on factors outside the subject of the situation, the convenience, but not necessarily an obligation strictly speaking, but rather ethical, moral or professional” (\citealt[41]{Pila2022}, my translation). 


\ea \label{ex:roncero:25}
    Resian \citeauthor{Pila2022} (\citeyear[36]{Pila2022}, my translation and glosses).\\
    \gll Zütra ćon det no jabulko mimo sïno.\\
    tomorrow want-\textsc{1sg} give.\textsc{inf} one apple.\textsc{acc.sg} \textsc{poss.1sg.dat.sg} son.\textsc{dat.sg} \\
    \glt ‘Tomorrow I will (certainly, as planned) give an apple to my son.’
\z

\ea \label{ex:roncero:26}
    Resian \citeauthor{Pila2022} (\citeyear[36]{Pila2022}, my translation and glosses).\\
    \gll Zütra bon dala no jabulko mimo sïno.\\
    tomorrow \textsc{be.fut.1sg} give.\textsc{f.sg} one apple.\textsc{acc.sg} \textsc{poss.1sg.dat.sg} son.\textsc{dat.sg} \\
    \glt ‘Tomorrow I will (probably) give an apple to my son.’
\z

\ea \label{ex:roncero:27}
    Resian \citeauthor{Pila2022} (\citeyear[45]{Pila2022}, my translation and glosses).\\
    \gll Se bojïn spest. Ći bon spadel, du be ma me wzdïnot?\\
    \textsc{reflex}.\textsc{1sg} be\_afraid.\textsc{1sg} fall.\textsc{inf} if \textsc{be.fut.1sg} fall.\textsc{m.sg} who then have.\textsc{3sg} \textsc{1sg.acc} lift \\
    \glt ‘I'm afraid of falling. If I fall, who will pick me up (lit. who has to carry me)?’
\z

\ea \label{ex:roncero:28}
    Resian \citeauthor{Pila2022} (\citeyear[41]{Pila2022}, my translation and glosses).\\
    \gll Ko bon wbïla to kököš, jä man si skühat nö malö bröda!\\
    when \textsc{be.fut.1sg} kill \textsc{art} chicken \textsc{1sg.nom} have.\textsc{1sg} \textsc{reflex} cook one little broth.\textsc{gen.sg} \\
    \glt ‘When I (will) kill the hen, I will make myself (lit. I have to cook) some broth.’
\z

\subsubsection{Within European de-volitives} \label{sec:roncero:3.2.2}
Outside of the Balkan context, \citet[322]{Dahl2000b} finds a second de-volitive future construction family within the Germanic languages of the North Sea: Danish, Faroese and Norwegian Bokmål (\textit{vil}); Frisian (\textit{wal}); and perhaps also Yiddish (\textit{vel}). The English future tense construction \textit{will} is also originally a volitive verb.\footnote{According to \textit{Oxford English Dictionary} (\url{https://oed.com}, last checked 20/03/2023) in Contemporary English the verb \textit{will }can mean “[t]o wish, desire; sometimes with implication of intention” although it classifies it as obscure, rare and archaic.} According to \citet[322]{Dahl2000b}, it is the most grammaticalized future tense construction of all the languages in this construction family.

Other than the semantics of the auxiliary chosen as a base for the future tense, it is unlikely that there was an areal influence of the North Sea/Germanic \textit{will} (and its cognates) in the development of West Polesian \textit{xoʧu}. First, historically it is hard to prove (especially given that the closest Germanic language that could have been affected, High German, did not develop it). And second, there is no phonological correspondence between West Polesian \textit{xoʧu} and Germanic \textit{will} (and cognates); whilst \textit{xoʧu} can be clearly related to the South Slavic forms like \textit{(hó)ču}.

Thus, where does the West Polesian de-volitive future stand with regards to the Balkan (South Slavic) de-volitive? On the one hand, we can hypothesize that the West Polesian de-volitive is a loan or calque of South Slavic. There are a few other features in West Polesian which look like South Slavic (mostly Serbian; e.g. vocative masculines in \textit{–u}; or the de-obligative futures).\footnote{In fact, some speakers have told me that there was a wave of migration from the south (mostly around contemporary Serbia) to the lands of West Polesie around the XV century (which would be reflected in the common surnames). However, I am unaware of any scientific sources on this topic.} But on the other hand, all the evidence collected so far shows that the de-volitive future in West Polesian is an independent development departing from a common base that has only grammaticalized in a few modern Slavic varieties, as we have proposed for the de-obligative. According to \citet[162]{Danylenko2011}, in Common Slavic “the verbs \textit{iměti }‘to have’ and \textit{xotěti / xъtěti} ‘to want’ were used in such P[eriphrastic] C[construction]s as de-modal extensions rather than auxiliaries.” In any case, the final output of the contemporary semantics of \textit{xoʧu }diverges from the rest of the construction families (typological expectations).

\begin{figure}
    \centering
    \includegraphics[width=0.8\linewidth]{figures/roncero/figure2.pdf}
    \caption{Distribution of de-volitive future in Europe. Sources (\cite{Dahl2000b}, \cite{Danylenko2011}, \cite{Duchet1995}, \cite{Gvozdanovic1995}, \cite{Kuteva1995}, \cite{Mackevic1959}). (CC-BY-SA Kristian Roncero).}
    \label{fig:roncero:2}
\end{figure}

\section{Summary and discussion} \label{sec:roncero:4}
We have seen that West Polesian has several future tense constructions, some of which are further or less advanced in their grammaticalization process. For most of them, especially for the two under discussion here, they seem own creations from available materials in Slavic, as none of its neighbors (in Standard forms, except for Belarusian \textit{mecca}) have any of them synchronically (and most of the contact with them only started about 70 years ago). 

Against the general typological expectation (in \tabref{tab:roncero:1} represented by English), the West Polesian de-volitive and de-obligative constructions have both retained some modal value, but they have inverted their original semantic functions. South Slavic (particularly Bulgarian and Macedonian) retains both (cognate) constructions, but it has generalized their use; that is to say, it has gone further in grammaticalization (especially if it is understood as ‘semantic bleaching’) and Resian Slavic keeps three constructions, which have developed fairly predictable nuances.



\begin{table}
\centering
\begin{tabularx}{\textwidth}{QQp{2.3cm}QQ} %{p{2cm} p{1.5cm} X p{1.5cm} X}
\lsptoprule
                    & {English}     & {South Slavic} & {Resian}            & {W. Polesian} \\
  \midrule
  Intention, Desire & \textit{will} & (BCMS)\newline \textit{(hó)ču} ‘want’;  & \textit{ćon} ‘want’ & \textit{majujsj(a)} ‘have’\\
\cline{1-2} \cline{4-5}
  Prediction        & \textit{will} & (BG/MKD)\newline +\textit{šte/ḱe} ‘want’\newline -\textit{[nja]ma} ‘have \newline [not]’ & \textit{bon} ‘be’   & (copular, synthetic, perfective, or become futures)\\
\cline{1-2} \cline{4-5}
  Obligation        & \textit{shall}&                & \textit{man} ‘have’ & \textit{xoʧu} ‘want’\\
  \lspbottomrule
\end{tabularx}
\caption{Examples of uses of future tense constructions}
\label{tab:roncero:1}
\end{table}


How did West Polesian end up in this situation? In other words, how did both constructions end up inverting their original meanings? Given then scarcity of data and the impossibility to carry out more fieldwork, I shall be cautious about any conclusion. Nevertheless, data seem to suggest that this overlapping of meanings is merely accidental, i.e. both forms were in the process of desemanticization/generalization, and they both ended up finding ‘an empty nest’ at different times in history. 

According to Spencer (p.c.), English also had a similar type of ‘inversion’. The de-volitive \textbf{\textit{will}} used to mean plain volition (as in \textit{as you will}), which then became a future auxiliary. Conversely, \textbf{\textit{want}} used to mean necessity or need (as in \textit{for they shall want for nothing}), but then filled the slot vacated by \textbf{\textit{will}}.

From a more traditional approach, this semantic ‘dissonance’ can be also the result of a saturated inventory of future tense constructions in West Polesian, which may have ‘forced’ certain members to specialize over time, which would be in line with \citegen{Aronoff2016} theory of Morphological Niches. This may explain what happened in Resian. Nevertheless, based on other observations about other morphosyntactic parameters in West Polesian, this seems to contradict the overall picture. I have gathered fine grained evidence of the same speakers using multiple endings (\cite{roncero2021revisiting}; \cite{Roncero2022}) for the same slot, without having any additional semantic, discursive or even sociolinguistic nuance between them (\cite{BošnjakBoticaHržica2016}; \cite{Thornton2011}); yet conditioned by its sociolinguistic setting (cf. \cite{Dorian2010}). I shall admit though, that in some (not all) cases, language contact seemed to be one of the sources of overabundance; whereas, the neighboring varieties seem to have very limited or no effect on the development and re-semantization of these two future tense constructions. I shall add to this that, although marginally, the reflexive auxiliary \textit{mecca} was already attested in Middle Belarusian, which \textit{mutatis mutanda}, may have already the intentional nuance. As for the de-volitive, the fact that we also find traces of it in Rusyn, suggests it existed in earlier stages of East Slavic but it was not fully grammaticalized. Yet, it is still a mystery at which point in history did West Polesian de-volitive turn into de-obligative. In sum, the evolution of the de-volitive and the de-obligative future constructions in West Polesian opens further questions about grammaticalization and semantics, as well as typology and sociolinguistics.

In any case, despite the overabundance of forms, in a big part, also conditioned by language contact, the semantics seem to follow a predictable pattern. By contrast, synchronically the West Polesian de-volitive form has ended up expressing necessity, whereas the de-obligative is used to indicate intention and desire. 



\section*{Abbreviations}
Glossing abbreviations in this paper are based on the \textit{Leipzig Glossing Rules} (2015 revision), with some additions.\vspace{10pt}

\begin{tabularx}{.5\textwidth}{@{}lQ@{}}
\textsc{1, 2, 3} & first, second, third person (respectively) \\
\textsc{abs} & absolutive \\
\textsc{acc} & accusative \\
\textsc{bcms} & Bosnian-Croatian-Montenegrin-Serbian \\
\textsc{bel} & Standard (Literary) Belarusian \\
\textsc{dat} & dative \\
\textsc{det} & determiner \\
\textsc{dem} & demonstrative \\
\textsc{erg} & ergative \\
\textsc{f} & feminine \\
\textsc{fut} & future \\
\textsc{gen} & genitive \\
\textsc{gradnm} & greater adnumerative \\
\textsc{imp} & imperative \\
\end{tabularx}%
\begin{tabularx}{.5\textwidth}{@{}lQ@{}}
\textsc{inf} & infinitive \\
\textsc{ins} & instrumental \\
\textsc{loc} & locative \\
\textsc{m} & masculine \\
\textsc{n} & neuter \\
\textsc{neg} & negation \\
\textsc{nom} & nominative \\
\textsc{pl} & plural \\
\textsc{poss} & possessor \\
\textsc{prf} & perfective \\
\textsc{ptcp} & participle \\
\textsc{pst} & past \\
\textsc{q} & question particle/marker \\
\textsc{refl} & reflexive \\
\textsc{res} & Resian Slavic \\
\textsc{sg} & singular \\
\textsc{voc} & vocative \\
\end{tabularx}

\section*{Acknowledgements}
This work has only been possible thanks to the participation of the West Polesian community in Belarus. Endless thanks for their patience and help. Many thanks as well to the Belarusian Charity (UK), The Gilchrist Educational Trust (UK) and The Philological Society (UK) who funded part of the fieldtrip. I shall acknowledge the support of the UK Arts and Humanities Research Council (grant no. AH/T002859/1), which funded part of the research time that allowed me to summarize my results into this chapter. Finally, thanks to Erik Tauschel, my Student Assistant, for helping with the formatting.

\begin{paperappendix}

\section{List of participants} \label{sec:roncero:app1}
I am providing the names of the participants who chose to have their identities disclosed. To all of those who show it here and for all the others who chose to remain anonymous, endless thanks.\vspace{10pt}

\noindent \textbf{Imianin [Torokan] (Drahičyn)}\\
\begin{tabularx}{.5\textwidth}{@{}lQ@{}}
T2 & Mazuka, Mikalaj A. \\
T5 & Mazuka, Kacjaryna I. \\
T6 & Barysjuk, Maryja I. \\
T7 & Marusja \\
T9 & Hardzejčuk, Nadzeja I. \\
T12 & Sakaloŭska, Maryja P. \\
Tor 1 & Baba Ženja \\
\end{tabularx}
\vspace{10pt}

\noindent \textbf{Pare (Pinsk)}\\
\begin{tabularx}{.5\textwidth}{@{}lQ@{}}
P2 & Maryja Uladzmiraŭna \\
\end{tabularx}
\vspace{10pt}

\noindent \textbf{Žydča (Pinsk)}\\
\begin{tabularx}{.5\textwidth}{@{}lQ@{}}
Z2 & Haŭrylčuk, Nina \\
Z10 & Žydečkaja, Sjarafima I. \\
\end{tabularx}
\vspace{10pt}

\noindent \textbf{Haloŭčyci, Tolkovo (Drahičyn)}\\
\begin{tabularx}{.5\textwidth}{@{}lQ@{}}
TL1 & Kanavalčyk, Vera I. \\
HL1 & Mancèvič, Dzina I. \\
\end{tabularx}

\end{paperappendix}
\sloppy
\printbibliography[heading=subbibliography,notkeyword=this]
\end{document}



